We overview previous work on calculating and applying meta-gradients.

\subsection{Calculating metagradients}
Previous work estimates the metagradient for large-scale models via one of
two broad families of methods: implicit differentiation and automatic (explicit)
differentiation. Note that in previous literature, synonyms for metagradient
include ``hyper-gradient'' and ``outer gradient.''

\paragraph{Implicit differentiation.}
One family of methods aims to {\em approximate} the metagradient.
To illustrate the idea behind such approaches, suppose that the
learning algorithm
$\mathcal{A}$ returns a model state $\theta$ that minimizes a strongly
convex loss function $\mathcal{L}(z, \theta)$. Here, the implicit function
theorem tells us that
\begin{align}
  \label{eq:ift}
  \nabla_z f(z) = \overbrace{
    \left(\frac{d\phi}{d\theta}\bigg|_{\theta = \mathcal{A}(z)}\right)
  }^{\substack{1 \times p \text{ gradient of output} \\ \text{wrt.
  final params}}}
  \underbrace{
    \left(
      \frac{\partial^2 \mathcal{L}(z, \theta)}{\partial
      \theta^2}\bigg|_{\theta = \mathcal{A}(z)}
    \right)^{-1}
  }_{\substack{p \times p \text{ inverse Hessian of loss}
  \\ \text{wrt. final params}}}
  \overbrace{
    \left(\frac{\partial^2 \mathcal{L}(z, \theta)}{\partial
    \theta\,\partial z}\bigg|_{\theta = \mathcal{A}(z)}\right).
  }^{
    \substack{p \times n \text{ Jacobian of loss gradient}
    \\ \text{wrt. metaparameters}}
  }
\end{align}
The form of \eqref{eq:ift} yields efficient and accurate estimators for
metagradients of models learned by minimizing a strongly convex
loss~\citep{bertrand2020implicit,bertrand2022implicit,kolter2020deep,blondel2022efficient,scieur2022curse}.
Such approaches can extend to estimate metagradients of large-scale, non-convex
learning
algorithms~\citep{bengio2000gradient,koh2017understanding,rajeswaran2019meta,finn2017model,lorraine2020optimizing,chen2020stabilizing,bae2022if},
but lose any correctness guarantees. Indeed, applying this class of methods in
large-scale settings is challenging as doing so requires (a) assuming conditions
on the learning algorithm (e.g., Hessian invertibility, continuous
differentiability) and (b) efficiently approximating the inverse Hessian (in
practice, typically at the cost of estimate accuracy). Finally, implicit
function-based approaches are fundamentally limited in that they can only
differentiate with respect to metaparameters expressed in the loss function
(e.g., these methods can differentiate with respect to the weight decay, but not
learning rate).

\paragraph{Automatic (explicit) differentiation.} Beyond implicit
differentiation approaches, there is a long line of work on directly calculating
metagradients with AD (see Section~\ref{sec:computing}). Previous work has used
AD to estimate metagradients of learning algorithms ranging from those with
convex objectives to small neural
networks~\citep{hara2019data,maclaurin2015gradient,franceschi2017forward,micaelli2021gradient,zhang2021idarts,chandra2022gradient,scieur2022curse}.
As detailed in Section~\ref{sec:computing}, the primary challenge with
(reverse-mode) AD-based approaches to meta-differentiation is storing the
intermediate products required for the backward pass. To circumvent this
challenge, previous work either (a) only considers settings that are small
enough that is possible to differentiate while requiring space that is linear in
the number of iterations (i.e., 2 layer networks on MNIST), (b) uses
forward-mode
AD~\citep{franceschi2017forward,micaelli2021gradient,chandra2022gradient} (which
  requires no extra storage at the cost of additional compute that
  scales linearly
with metaparameter dimension), (c) only \textit{approximates} the metagradient
by calculating over only a few training
steps~\citep{liu2018darts,chen2020stabilizing,finn2017model}, or uses (d) a
reversible learning algorithm~\citep{maclaurin2015gradient}. The fourth category
is a promising direction for reducing space requirements when computing
large-scale metagradients, but current approaches require (a) representing model
parameters in a fixed-precision format (which current large-scale learning
algorithms do not support) in addition to restricting the algorithm to be
reversible (e.g., SGD and standard GD do not qualify). A common thread is that
algorithms computing metagradients with AD often suffer from numerical
instability and overflow issues~\citep{micaelli2021gradient,scieur2022curse}. In
relation to previous work on AD, \replay{} (Section \ref{sec:computing})
can be seen as a strategy for choosing gradient
checkpointing~\citep{chaitin1981register,briggs1992rematerialization,zweig2000exact,griewank2008evaluating,chen2016training}
locations in the compute graph (an NP-complete task in
general~\citep{naumann2008optimal}).

\subsection{Applying metagradients}
Previous work applies metagradients to optimize training setup, including
distillation~\citep{maclaurin2015gradient,lorraine2020optimizing}, training data
selection~\citep{hara2019data,engstrom2024dsdm},
meta-learning~\citep{finn2017model,rajeswaran2019meta,hospedales2021meta},
learning rate/weight decay
selection~\citep{micaelli2021gradient,chandra2022gradient}, tuning data
augmentation~\citep{lorraine2020optimizing}, and architecture
search~\citep{maclaurin2015gradient,liu2018darts,zhang2021idarts}. Beyond
optimizing metagradients, methods in data attribution apply metagradients to
(Taylor) estimate the effect of dropping training data on model
predictions~\citep{koh2017understanding,grosse2023studying,park2023trak}. To the
Previous works either (a) calculate metagradients directly with AD (made
feasible by working in a very small-scale learning setting) or (b) estimate the
metagradient with an implicit function-based approach.

